\section{Introduction}
\quad In recent decades, advances in experimental techniques and data acquisition have made it possible to record high-dimensional time-series data in a wide range of biological, physical, and engineered systems. Many such interacting systems can be described as noisy network dynamics \cite{strogatz2001exploring, Mark2003, Mark2018}, in which heterogeneous components interact through weighted and possibly directed connections while being subject to stochastic fluctuations. These systems naturally generate time-series data from which covariance can be measured directly. From a physical perspective, many living systems are maintained through continuous exchanges of energy, matter, and information with their surroundings and therefore typically operate far from thermodynamic equilibrium \cite{PhysRevLett.71.2401, 10.1143/PTPS.130.29, sekimoto2010stochastic, PhysRevE.96.032123, PhysRevE.96.012601, Shiraishi2023}. Examples include metabolic processes and molecular motors. In such nonequilibrium systems, energy is continuously consumed and entropy is produced, while detailed balance can be broken by asymmetric interactions or by heterogeneous noise sources acting on different components.\par

One important inverse problem in complex networks is network reconstruction from time-series data. In many situations, the detailed connectivity is not directly accessible, whereas the node dynamics can be measured. Previous studies have shown that fluctuations can encode structural information and that network connectivity can be inferred from network dynamics \cite{PhysRevLett.104.058701, ching2013, ching2015, ching2017, Lai2017, ChengLai2022}. These results highlight the close relation between fluctuations and interactions in networked systems. Reconstruction alone, however, does not answer a closely related and equally important question: once the parameters of a noisy network are changed, how can one predict the resulting response of the system from properties of its original steady state? In many situations, the microscopic dynamics may be complicated, yet the system remains close to a stable fixed point and its fluctuation statistics can still be measured. This motivates a response theory formulated directly for noisy network dynamics near stable fixed points, where one can ask how perturbations of node parameters, couplings, or noise strengths modify steady-state observables.\par

In statistical mechanics, linear response theory (LRT) provides a general framework for predicting a system's response to small perturbations from correlation functions measured in the unperturbed state \cite{RKubo_1966, 10.1063/5.0003135}. In equilibrium systems, this connection is expressed by the fluctuation-dissipation theorem (FDT). Many realistic networked systems, however, do not operate near thermodynamic equilibrium, but instead settle into nonequilibrium steady states (NESS). Although LRT has been extended to NESS in several theoretical frameworks \cite{Seifert_2010, Mehl_2010, PhysRevLett.117.180601, PhysRevLett.134.157101}, these approaches do not directly provide explicit matrix-level susceptibilities for large networked systems. Therefore, there is a need for a practical framework that can be computed directly from the unperturbed system and then tested against perturbation experiments or simulations. This limitation becomes even more evident when one wishes to describe not only the difference between two steady states, but also the transient relaxation induced by a time-dependent perturbation protocol. Moreover, in coupled dynamical systems, local perturbations generally do not remain local. Perturbations of node parameters or interaction strengths can modify fluctuations and correlations throughout the system. In practical terms, local changes such as node removal, link failure, or targeted interventions can induce nontrivial changes in global observables \cite{NetworkDisruption}. A predictive framework for such responses is therefore useful not only for understanding noisy interacting systems, but also for assessing the effects of controlled perturbations and disruptions.

In this work, we address this problem for network dynamics fluctuating near stable fixed points. Each node is driven by Gaussian white noise and interacts with other nodes through weighted directed couplings, while also possessing its own intrinsic dynamics. Linearizing the dynamics around the noise-free stationary state leads to a Jacobian matrix $\mathbf{Q}$ and a steady-state equal-time covariance matrix $\mathbf{K}_0 \equiv \langle \delta \vec{x}\,\delta \vec{x}^{\top} \rangle$ satisfying a Lyapunov equation. Within this linearized Ornstein-Uhlenbeck description, the steady-state distribution is known exactly as a multivariate Gaussian. This exact solvability allows the first-order response functions to be derived explicitly from unperturbed steady-state quantities. For the original nonlinear noisy network, the same formulas should therefore be understood as the corresponding local Gaussian approximation near the stable fixed point. Within this scope, we develop a linear-response framework for both steady-state and time-dependent perturbations. Our central focus is the response of $\mathbf{K}_0$, for which we derive explicit and computable formulas under perturbations of intrinsic node parameters, connection weights, and noise strengths. The framework applies to both equilibrium steady states and nonequilibrium steady states of the linearized dynamics, where detailed balance is generally broken by asymmetric couplings or heterogeneous noise amplitudes. In this sense, the contribution of this work is not a general reformulation of NESS response theory for arbitrary nonlinear stochastic systems, but an operational response framework for linearized noisy networks near fixed points. The responses of means and forces further show that the same formalism also captures transient network relaxation.\par

The paper is organized as follows. In Sec. II, we introduce the noisy network model and its linearization near a stable fixed point. In Sec. III, we establish a linear response framework for both equilibrium and NESS in the linearized Langevin dynamics. In Sec. IV, we derive the steady-state linear response of $\mathbf{K}_0$ for our noisy network model. In Sec. V, we compare the theoretical predictions with Langevin simulations for both equilibrium and NESS cases. In Sec. VI, we extend the formalism to time-dependent perturbations and discuss the corresponding dynamical response.
